\documentclass[a4j,12pt,dvipdfmx]{jsreport}

% 数式
\usepackage{amsmath,amsfonts}
\usepackage{bm}

% 画像
\usepackage[dvipdfmx,top=35truemm,bottom=30truemm,left=30truemm,right=30truemm]{geometry}
\usepackage[dvipdfmx]{graphicx}
% biblatex
\usepackage[style=numeric, sorting=none, maxbibnames=1000]{biblatex}
\addbibresource{graduation_thesis.bib}

% itembox
\usepackage{ascmac}
% breakitemboxのためのpackage
\usepackage{itembkbx}
% \usepackage{tcolorbox} 画像が見えなくなってしまう原因
% \usepackage{framed} 枠

% フォント
\usepackage[ipaex]{pxchfon}
\listfiles
\usepackage{float}
\usepackage{multicol}
\usepackage{multirow}
\usepackage{paralist}
\usepackage{caption}
% \captionsetup{format=plain, justification=centering}
\usepackage{setspace}
\usepackage{algorithm}
\usepackage{algpseudocode}
\usepackage[dvipdfmx,hidelinks]{hyperref}
\usepackage{pxjahyper}
\usepackage{fancyhdr}
\usepackage{ltablex}
\usepackage{capt-of}
\usepackage{tabularx}
\usepackage[export]{adjustbox}
\usepackage{booktabs}
\usepackage{enumitem} % itemizeの行間 これは下になければならない
\usepackage{threeparttable} % 表の脚注
\usepackage{titlesec}
\usepackage{makecell}
\usepackage{forest}
\usepackage{fancybox}
% titlesec と \paragraph の相性問題を回避（run-in 見出しにする）
\titleformat{\paragraph}[runin]
  {\normalfont\normalsize\bfseries}
  {}{0pt}{} % ←ここには文字や \noindent 等を入れない
\titlespacing*{\paragraph}
  {0pt}{1.0ex}{0.8em}
\titleformat{\chapter}[hang]{\LARGE\bfseries}{\prechaptername\thechapter\postchaptername}{1em}{}

\setlength{\headheight}{16.5pt}
\captionsetup{format=hang}
\captionsetup[figure]{font=small}
\captionsetup[table]{font=small}
\begin{document}

% ▼▼▼ 仕掛け：main じゃないとき（単体のとき）だけヘッダー設定 ▼▼▼
\ifdefined\ismain\else
   \pagestyle{fancy}
   \lhead{\rightmark}
   \renewcommand{\chaptermark}[1]{\markboth{第\ \thechapter\ 章\ #1}{}}
   \setlength{\headheight}{16.5pt}
   \titleformat{\chapter}[hang]{\LARGE\bfseries}{\prechaptername\thechapter\postchaptername}{1em}{}
\fi
% ▲▲▲▲▲▲▲▲▲▲▲▲▲▲▲▲▲▲▲▲▲▲▲▲▲▲▲▲▲▲▲▲▲▲▲▲

\chapter{結論}\label{chap:conclusion}

本研究では，
漫画制作工程において重要な役割を担う「ネーム」を，
作品の下書き画像そのものではなく，
物語文脈と演出意図を内包した
構造化された中間表現として捉え直し，
脚本に基づく漫画ネーム生成という新たな課題設定を提示した．
従来のレイアウト生成研究が，
主として機能的媒体や要素配置の最適化を対象としてきたのに対し，
本研究は，
物語性や演出上の意思決定と不可分な
漫画特有の配置問題を扱う点に特徴がある．

この課題に取り組むため，
本研究では，
大規模視覚言語モデル（LVLM）を活用し，
漫画画像に紐づく物語情報および演出情報を体系的に抽出・構造化した
Manga109ScriptDataset および Manga109NameDataset を構築した．
特に，
人手による読み順アノテーションと
独自の Visual Prompting 手法を組み合わせることで，
LVLM 単体では困難であった
漫画特有の物語構造や演出情報の記述を可能にし，
脚本およびネーム生成のための学習基盤を整備した．

さらに，
構築したデータセットを用いて，
脚本を入力とし，
演出情報およびレイアウト情報を生成する
漫画ネーム生成モデルを提案した．
大規模言語モデル(LLM)を基盤とした複数の生成設定を比較した結果，
脚本情報を明示的に入力として与えることで，
既存のレイアウト生成手法を上回る性能が得られることを
定量的・定性的に示した．
また，
演出情報をどの段階で統合するかという設計の違いなどを複数検討し設計指針を整理した．

一方で，
演出情報は本質的に多義的であり，
単一の正解に基づく評価が難しいという課題も明らかとなった．
本研究では，
GNN を用いた間接評価や可視化分析を通じて，
演出情報の品質を検討したが，
明らかに文脈と乖離した演出や，
制作上不適切な生成結果を抑制するためには，
演出情報の妥当性や一貫性を
より定量的に捉える評価指標の設計が今後必要である．

脚本を条件としてネームを生成するという枠組みは，
最終成果物を直接生成するのではなく，
制作工程の中核に位置する意思決定を支援するという点で，
創作者の創造性を尊重しつつ，
制作効率の向上に寄与し得る可能性をもつと考えられる．
本研究で示したデータセット，モデル設計，および評価結果は，
脚本に基づく漫画ネーム生成という新たな研究領域における
基礎的なベースラインとなるものであり，
今後の発展的研究や実制作支援への応用に向けた
重要な足がかりとなることを期待する．


% ▼▼▼ 仕掛け：main じゃないとき（単体のとき）だけ参考文献を表示 ▼▼▼
\ifdefined\ismain\else
   \printbibliography[title=参考文献]
\fi
% ▲▲▲▲▲▲▲▲▲▲▲▲▲▲▲▲▲▲▲▲▲▲▲▲▲▲▲▲▲▲▲▲▲▲▲▲

\end{document}